\documentclass[11pt]{article}
\usepackage[margin=1in]{geometry}
\usepackage{amsmath,amssymb,booktabs,graphicx,hyperref,microtype,float}
\newfloat{algorithm}{htbp}{loa}
\floatname{algorithm}{Algorithm}
\newcommand{\N}{\mathcal N}
\newcommand{\dd}{\mathrm d}
\newcommand{\E}{\mathbb E}
\newcommand{\rank}{\operatorname{rank}}
\title{A Ditlevsen--Karppinen transition-density competitor for the
Dacca cholera benchmark}
\author{Computational note}
\date{\today}

\begin{document}
\maketitle

\section{Question and evaluation target}

Reviewer~1 proposed constructing tractable Gaussian transitions using the
strong-order-1.5 approximation of Ditlevsen and Samson (DS), then using those
transitions in SAEM, Fisher-identity gradient estimation, PGAS, or a particle
Gibbs method stabilized by Karppinen, Singh, and Vihola (KSV).  The statistical
criterion in the DMOP paper is not the complete-path pseudo-likelihood.  It is
the log likelihood of the Euler-discretized POMP, independently evaluated by a
bootstrap particle filter.  Accordingly, a candidate method is successful only
if its estimated parameter vector has a competitive Euler-model log likelihood.

The paper used 100 searches from the same wide box, compared methods under
similar elapsed time, evaluated selected parameter vectors with 5,000 particles
and 36 replications, and emphasized both the right tail and median of the final
log likelihoods.  Its best comparable-effort result is $-3744.17$ for
IFAD-$0.97$; extended IFAD-$0.97$ reaches $-3744.19$ with median $-3744.22$.

\section{Dacca state representation}

For the benchmark, $c=\alpha=1$, $\rho=Y_0=0$, and $\delta=0.02$ are fixed.
Thus the inapparent-infection coordinate remains zero.  In population-fraction
coordinates, the active state is
\[
 X=(S,I,R^1,R^2,R^3,M)^\mathsf T\in\mathbb R^6,
\]
where $M$ accumulates disease deaths within a month and is reset after each
measurement.  With drift $b_\theta(X,t)$, the single diffusion vector field is
\[
 g_\theta(X,t)=\sigma S I\,(-1,1,0,0,0,0)^\mathsf T,
 \qquad \dd X_t=b_\theta(X_t,t)\dd t+g_\theta(X_t,t)\dd B_t.
\]
The observation model is exactly the manuscript model,
$Y_n\mid X_{t_n}\sim\N(P_{t_n}M_{t_n},
\tau^2P_{t_n}^2M_{t_n}^2)$.

\section{Diffusion rank, transition rank, and bracket depth}

It is essential not to confuse the two ranks.  Hypoellipticity means that the
\emph{instantaneous diffusion matrix is not full rank} while the finite-time
transition nevertheless has a smooth density.  DS assume
$V\in\mathbb R$, $U\in\mathbb R^p$, and
\[
 C(v,u)=\begin{bmatrix}0_{1\times p}\\\Gamma(v,u)\end{bmatrix},
 \qquad
 \Gamma(v,u)=\operatorname{diag}\{\sigma_1(v,u),\ldots,
 \sigma_p(v,u)\},\quad \sigma_j>0.
\]
Thus $C$ has rank $p<p+1$ by construction.  The directly forced rough block
already supplies $p$ directions and the first drift bracket supplies the one
missing smooth direction.  The order-1.5 term is exactly what gives their
leading approximate covariance a nonzero variance of order $\Delta^3$ in that
smooth coordinate for sufficiently small $\Delta$.

Equivalently, the DS instantaneous diffusion has rank $d-1$ in a $d$-state
system.  Dacca has $d=6$ active coordinates but diffusion rank one, hence
codimension five rather than one.  Diffusion rank is invariant under an
invertible change of state coordinates, so state dependence of the coefficient
$\sigma SI$ does not put Dacca into the DS model class.

Dacca is potentially hypoelliptic too, but it is not in this depth-one form.
Its active state has dimension six and its instantaneous diffusion has rank one
away from $SI=0$ and rank zero on that boundary.  At the manuscript parameter
and state point, the normalized columns formed with the
Stratonovich drift $b^{\mathrm S}=b-\tfrac12(Dg)g$,
\[
 g,\ [b^{\mathrm S},g],\ [b^{\mathrm S},[b^{\mathrm S},g]],\ldots
\]
gain numerical ranks $1,2,3,4,5,6$ only through bracket depth five.  This is
evidence along one repeated drift-bracket chain that the true finite-time SDE
transition can be full-dimensional; it is neither a proof of global
hypoellipticity nor an argument that the diffusion should be full rank.

The difficulty is instead that a literal first-bracket generalization of the
DS leading covariance does not contain all those directions.  Let
$J_b=\partial b_\theta/\partial x$.  Locally freezing the state-dependent
diffusion gives

\begin{align}
 Q^{(1.5)}_h(x)
 &=hgg^\mathsf T
 +\frac{h^2}{2}\{(J_bg)g^\mathsf T+g(J_bg)^\mathsf T\}
 +\frac{h^3}{3}(J_bg)(J_bg)^\mathsf T\\
 &=A_2H_2A_2^\mathsf T,\\
 A_2&=\begin{bmatrix}h^{1/2}g&h^{3/2}J_bg\end{bmatrix},
 & (H_2)_{ij}&=\frac{1}{i+j+1},\quad i,j\in\{0,1\}.
\end{align}
Therefore
\[
 \rank\{Q^{(1.5)}_h(x)\}\leq 2<6.
\]
This conclusion is not invalidated by the state dependence highlighted by the
reviewer.  Here $g(x)=q(x)e$ with the constant direction
$e=(-1,1,0,0,0,0)^\mathsf T$, so $(Dg)w=e\{(Dq)w\}$ is collinear with $g$ for
every $w$.  The Milstein and other diffusion-derivative terms in the full
strong-order-1.5 expansion therefore remain on the directly forced line; the
first drift-propagation term adds only the $J_bg$ direction.  The paper gives
its nondegenerate Gaussian construction for the $1+p$ structure with $p$
directly forced directions, not a six-state/one-driver, depth-five analogue.
This rank-two statement concerns the approximate \emph{finite-step
covariance}, not the rank-one instantaneous diffusion.  It shows that the DS
nondegeneracy argument for their $1+p$ depth-one model cannot simply be copied
to this longer-chain system.  Further stochastic Taylor/bracket terms, or a
block construction that accumulates them, are required.  The implementation
also uses the DS second-order mean
\[
 m_{\theta,h}(x)=x+h b_\theta(x)+\frac{h^2}{2}\mathcal L_\theta
 b_\theta(x),
\]
where $\mathcal L_\theta$ is the It\^o generator.

Adding a nugget makes numerical evaluation possible but changes the statistical
model, with the nugget dominating directions that the first-bracket expansion
leaves exactly deterministic.  Results labeled ``DS-local+nugget QML'' disclose
this change and are not presented as a literal execution of DS Algorithm~2.

\section{Which time scale is being compared}

One Euler step per measurement interval is indeed the measurement time scale.
The important distinction is between transition \emph{schemes} on that same
grid.  With $r$ substeps per month,
\[
 r=1:\ \text{one monthly DS transition},\qquad
 r=20:\ \text{twenty DS transitions of length $1/20$ month}.
\]
DS demonstrate inference with their discretization points at the observation
times.  Introducing $r-1$ intermediate latent states per month is an extension.
In every case below, goodness is assessed against the \emph{same} numerical
model used by the DMOP manuscript: the Euler POMP with 20 substeps per month.
A one-step Euler likelihood is therefore neither the target nor evidence about
the one-step DS fit.

Grid equality also says nothing about numerical stability.  The published
Dacca point has a three-stage immunity rate $3\epsilon=57.3$ per year.  Applied
to a decay mode with rate $\lambda$, the DS second-order conditional mean has
stability polynomial $R(z)=1-z+z^2/2$, $z=h\lambda$, whose real-axis
absolute-stability interval is $0\leq z\leq2$.  The monthly, half-monthly, and
five-substep values of $h(3\epsilon)$ are respectively $4.78$, $2.39$, and
$0.96$.  Thus failure at $r=1,2$ and viability beginning at $r=5$ are the
expected coarse-step behavior of this mean, not evidence that the measurement
and one-step grids were different.

\section{Higher-bracket completion as a diagnostic}

To test the strongest nearby density-based idea, define the local
controllability expansion
\begin{align}
 A_q(x,h)&=\left[
 h^{1/2}g,\frac{h^{3/2}}{1!}J_bg,\ldots,
 \frac{h^{q-1/2}}{(q-1)!}J_b^{q-1}g\right],\\
 (H_q)_{ij}&=\frac{1}{i+j+1},\qquad
 Q_h^{(q)}(x)=A_qH_qA_q^\mathsf T.
\end{align}
This is the covariance obtained by integrating the Taylor expansion of
$\exp(J_bs)g$ over a step.  For $q=2$ it reproduces the DS leading covariance;
for $q=6$ the unscaled Dacca local controllability matrix is generically full
rank.  The mean remains the DS second-order generator expansion.  This
``DS--Krylov'' completion is a new local-linear Gaussian approximation, not a
claim about the scope of the DS paper.  At fine $h$, its smallest directions
scale as high odd powers of $h$, so the transition becomes extremely
ill-conditioned even when it is algebraically full rank.  It is implemented as
a rank diagnostic and optional sensitivity analysis.  The primary executable
competitor instead uses the literal two-column DS leading covariance plus a
reported relative numerical nugget; neither choice is presented as an exact
application of DS to Dacca.

The earlier extended-Kalman quasi-likelihood optimizer has been withdrawn and
is not used as a DS result.  The executable competitor now uses a bootstrap
particle filter at every parameter iteration, draws one path from the forward
genealogy, evaluates the complete Gaussian pseudo-likelihood on that path, and
takes a projected stochastic Fisher-score step.  The score average uses the DS
gain schedule.  Because the Dacca complete pseudo-likelihood has not been
shown to have DS's curved exponential-family form, there is no SAEM M-step in
the executable method: this is stochastic gradient ascent, not DS Algorithm~2.
Every candidate is judged afterward by the original Euler-20 particle filter
rather than by its training pseudo-likelihood.

More precisely, write $\widetilde p_{\Delta,\theta}$ for the regularized DS
Gaussian pseudo-model with latent path $X$.  Fisher's identity gives
\[
 \nabla_\theta \widetilde\ell_\Delta(\theta;y)
 =\E_\theta\!\left[
   \nabla_\theta\log \widetilde p_{\Delta,\theta}(X,y)\mid y
 \right].
\]
At iteration $m$, forward SMC approximates
$\widetilde p_{\Delta,\theta_{m-1}}(X\mid y)$ and one complete genealogy
$X^{(m)}$ is drawn from its terminal weighted empirical measure.  Automatic
differentiation is then applied to
$\log \widetilde p_{\Delta,\theta}(X^{(m)},y)$ while holding $X^{(m)}$ fixed.
Thus
\[
 \widehat z_m=
 \left.\nabla_\theta\log \widetilde p_{\Delta,\theta}(X^{(m)},y)
 \right|_{\theta=\theta_{m-1}}
\]
is a one-path SMC approximation of the pseudo-model observed score.  It is not
an estimate of the original Dacca SDE/Euler score: the latter transition
density is unavailable, and the first-bracket Dacca covariance is singular
until the disclosed nugget is added.  No derivative is taken through particle
simulation, resampling, or ancestor selection.  Updating
$s_m=(1-a_m)s_{m-1}+a_m\widehat z_m$ and stepping along $s_m$ is a
Robbins--Monro score method.  In contrast, proper SAEM updates a
parameter-independent sufficient statistic $S(X^{(m)},y)$ and maximizes its
surrogate exactly (or with a justified generalized-EM step).

\section{Why KSV is relevant, and why it is not a direct patch}

At $r=1$ there are no artificial within-month latent points, so KSV does not
solve the monthly Taylor-approximation or rank-depth issues.  It can become
relevant when DS is extended to $r>1$.  In ordinary particle backward sampling,
the ancestor weight contains
$p_h(x_{t+h}\mid x_t)$.  As $h$ decreases, this density concentrates and almost
all backward weight is assigned to the existing ancestor.  Increasing the
number of Euler steps between monthly observations therefore degrades path
mixing.  KSV replace unit-step backward sampling by bridge updates spanning a
fixed physical-time block.  Their method directly targets this failure mode.

There is an important algorithmic distinction.  DS Algorithm~2 draws a fresh
path from an ordinary forward SMC approximation at every SAEM iteration; it
does not require a conditional-particle-filter Markov kernel.  KSV specifically
stabilize the backward update of a CPF/particle-Gibbs kernel.  Consequently KSV
is not required merely to run the DS-style forward-SMC competitor with twenty
substeps.  It would be relevant if that independent imputation were replaced
by particle Gibbs or another CPF path update.

KSV Assumption~7 requires evaluation of multi-step densities
$M_{u\mid\ell}(x_u\mid x_\ell)$ and simulation from bridge conditionals
$M_k(x_k\mid x_{k-1},x_u)$.  The nonlinear, state-dependent DS local transition
does not supply these objects.  A valid combination must introduce a tractable
time-varying linear-Gaussian proposal $M$ and put the density correction into
the Feynman--Kac potentials.  KSV also does not repair the rank-two DS
transition, and Dacca's complete likelihood is not the curved exponential
family required for DS's closed-form SAEM M-step.  One must either derive a
valid generalized-EM surrogate or use a separately labeled score method.

\begin{algorithm}[ht]
\caption{Full DS--Krylov--KSV stochastic-approximation score ascent
(design, not the executed baseline)}
\label{alg:combined}
\textbf{Input:} observations $y_{1:N}$; initial $\theta_0$; Euler step $h$;
particles $J$; SA gains $a_m$; physical-time block endpoints $\mathcal T$.
\begin{enumerate}
\item Construct the higher-bracket completed transition
\[
p^{K}_{\theta,h}(x_k\mid x_{k-1})=
\N\!\left(m_{\theta,h}(x_{k-1}),Q^{(6)}_{\theta,h}(x_{k-1})\right).
\]
\item At the current reference path, construct a globally time-varying
linear-Gaussian proposal $M_{\theta,1:T}$ whose multi-step marginals and bridge
conditionals are available by Kalman recursions.
\item Define correction potentials
\[
 G_{\theta,k}(x_{k-1},x_k)=
 \frac{p^{K}_{\theta,h}(x_k\mid x_{k-1})}
      {M_{\theta,k}(x_k\mid x_{k-1})}
 \times g_\theta(y_n\mid x_k)^{\mathbf 1\{k=t_n\}}.
\]
\item For $m=1,2,\ldots$:
  \begin{enumerate}
  \item Run a conditional particle filter for
  $(M_{\theta_{m-1}},G_{\theta_{m-1}})$.
  \item Apply KSV bridge backward sampling over $\mathcal T$, proceeding from
  the final block to the first and using the Gaussian bridge conditionals of
  $M$, to obtain $X^{(m)}_{0:T}$.
  \item Compute the complete-data score
  $z_m=\nabla_\theta\{\log p_\theta(X^{(m)}_0)+
  \sum_k\log p^K_{\theta,h}(X^{(m)}_k\mid X^{(m)}_{k-1})+
  \sum_n\log g_\theta(y_n\mid X^{(m)}_{t_n})\}$.
  \item Update the stochastic approximation
  $s_m=(1-a_m)s_{m-1}+a_mz_m$.
  \item Take a constrained stochastic-score parameter step, e.g.
  $\theta_m=\Pi_\Theta\{\theta_{m-1}+\eta_m\operatorname{Adam}(s_m)\}$.
  \end{enumerate}
\end{enumerate}
\textbf{Output:} $\widehat\theta$ and an independent Euler-particle-filter
estimate of $\ell_{\rm Euler}(\widehat\theta)$.
\end{algorithm}

Algorithm~\ref{alg:combined} is the minimal coherent combination of the two
papers for Dacca.  It is materially more than ``put KSV inside DS'': it requires
a higher-order rank completion, an auxiliary Gaussian bridge model, exact
importance correction for that proposal, a forward conditional-particle pool,
and a numerical M-step.

The code implements a narrower fixed-lower/fixed-upper blocked-interior
conditional SMC kernel derived from the bridge construction in KSV
Algorithm~8.  Its auxiliary bridge conditionals and $p^K/M$ corrections are
internally consistent, but it cannot reconnect the upper endpoint to a
different lower-boundary forward particle.  It is therefore neither full
Algorithm~8 nor CPF--BBS, and its fraction of changed interiors is not KSV's
probability of lower-boundary update (PLU).  Short Dacca checks found mixed
effects on parameter progress and no repair of the one-ancestor forward
genealogy.  It is excluded from the headline baseline unless a full forward
CPF--BBS implementation passes a separate invariance and mixing validation.

The executed numerical score method is a Fisher-identity Robbins--Monro
extension, rather than the closed-form curved-exponential-family SAEM used by
DS.  Neither it nor the blocked-interior diagnostic is used to claim that full
Algorithm~\ref{alg:combined} has been run.

\begin{thebibliography}{9}
\bibitem{ditlevsen19}
S. Ditlevsen and A. Samson (2019).
Hypoelliptic diffusions: filtering and inference from complete and partial
observations. \emph{Journal of the Royal Statistical Society, Series B},
81, 361--384.  \href{https://arxiv.org/abs/1707.04235}{arXiv:1707.04235}.

\bibitem{karppinen24}
S. Karppinen, S. S. Singh, and M. Vihola (2024).
Conditional particle filters with bridge backward sampling.
\emph{Journal of Computational and Graphical Statistics}, 33, 363--378.
\end{thebibliography}

\end{document}
